\chapter{Lanczos, Arnoldi, and Krylov--Schur}
\label{ch:krylovschur}

\chapterorient{The previous chapter posed the eigenproblem and handed its solve to
an opaque library: ``assemble the operators once and call \code{EPSSolve}.'' This
chapter opens that box. Inside is one of the most elegant constructions in numerical
linear algebra---a way to find a handful of a giant matrix's eigenvalues without ever
forming, or even looking at, the whole of it. The idea has three moving parts, and we
take them in order. First, \emph{Rayleigh--Ritz}: project the operator onto a small
Krylov subspace and read approximate eigenvalues off the tiny projected matrix; the
extreme eigenvalues come out first and best. Second, the \emph{restart problem}: the
subspace must stay small, but a naive restart throws away the good approximations you
worked to find. Third, \emph{thick-restart Krylov--Schur}: the elegant fix that keeps
the directions worth keeping and discards the rest, compressing the basis while
preserving its Krylov structure. We assemble these on top of the Arnoldi and Lanczos
recurrences of \cref{ch:orthog}, fit the inner shift-invert solve of
\cref{ch:eigenvalues} inside each step, and arrive at the full eigen-solve
stack---the exact loop that SLEPc and ARPACK run on Palace's behalf. The black box of
\cref{ch:eigenvalues} turns out to be code we could have written ourselves.}

% local helper: a Ritz value symbol used throughout (theta), already covered by
% standard math; no new macros needed. The basis matrix V_m uses \mat{V}.

\section{The setting: too big to factor, too many to want}
\label{sec:kssetting}

The eigenproblem of \cref{ch:eigenvalues} is the generalized system
\[
  \mat{K}\vx \;=\; \lambda\,\mat{M}\vx,
\]
with $\mat{K}$ and $\mat{M}$ the stiffness and mass operators assembled from the mesh.
Two facts about this problem shape everything that follows, and both are facts of
\emph{scale}. First, $\mat{K}$ and $\mat{M}$ are enormous---a fine mesh produces
millions of unknowns---and sparse, so we will never form a dense factorization of
either, let alone of the pair. We are allowed only to \emph{apply} the operators to
vectors (\cref{ch:matrixfree}) and to solve linear systems with them
(\cref{ch:multigrid}). Second, we do not \emph{want} all the eigenvalues. A cavity has
millions of resonant modes, but the engineer cares about the lowest dozen---the
fundamental and its first few overtones. We want a handful of eigenvalues near a target,
not the whole spectrum.

These two facts rule out the textbook eigenvalue algorithms. The QR
algorithm---the workhorse for small dense matrices---computes \emph{all} the
eigenvalues and needs the matrix in dense form; it is hopeless here on both counts. What
we need is a method that touches $\mat{K}$ and $\mat{M}$ only through matrix-vector
products and inner linear solves, and that spends its effort on the few eigenvalues we
asked for. That method is \emph{Krylov subspace projection}, and its practical form is
the subject of this chapter.

\begin{keyidea}
The eigenproblems of this book are large, sparse, and selective: millions of unknowns,
access only by operator-apply and linear-solve, but only a handful of wanted
eigenvalues near a target. This rules out dense methods that compute the whole spectrum.
The remedy is \emph{Rayleigh--Ritz projection onto a small Krylov subspace}: a small
projected eigenproblem yields the extremal eigenpairs, and \emph{thick-restart} keeps
the good ones while holding the subspace small. The entire algorithm uses $\mat{K}$ and
$\mat{M}$ only through the operator-apply and linear-solve we already have.
\end{keyidea}

The plan of attack inverts the difficulty exactly. We cannot work with the big operator,
so we project it onto a small subspace where we \emph{can} work with it---small enough
to hand to the dense QR algorithm. We cannot afford a huge subspace, so we keep it small
by restarting. And we want the eigenvalues near a target, not the extremes, so we feed
the iteration the shift-inverted operator of \cref{ch:eigenvalues}, which turns ``near
the target'' into ``extreme.'' Three subspace tricks, layered: projection, restart,
spectral transform. We build them in that order.

\section{Rayleigh--Ritz: a small projected eigenproblem}
\label{sec:rayleighritz}

Start with the first trick and the simplest version of it: no restart, no spectral
transform, just projection. Suppose we have built---by Arnoldi or Lanczos
(\cref{ch:orthog})---an orthonormal basis $\vq_1, \dots, \vq_m$ for the Krylov subspace
$\mathcal{K}_m(\mat{A}, \vv)$, stacked into the tall thin matrix
$\mat{V}_m = [\,\vq_1 \mid \cdots \mid \vq_m\,]$ with $\mat{V}_m^{\!\top}\mat{V}_m =
\mat{I}_m$. (For now read $\mat{A}$ as the operator whose eigenvalues we want; in the
eigenmode driver it will be the shift-inverted operator of \cref{sec:shiftinvert}.) The
question Rayleigh--Ritz answers is: \emph{what is the best approximation to $\mat{A}$'s
eigenpairs that we can extract from this subspace alone}?

\subsection{The projected operator}

We cannot diagonalize $\mat{A}$---it is too big. But we \emph{can} ask how $\mat{A}$
acts \emph{within} the subspace $\mathcal{K}_m$, by projecting its action down onto the
basis. The projected operator is the $m \times m$ matrix
\[
  \mat{T}_m \;=\; \mat{V}_m^{\!\top}\,\mat{A}\,\mat{V}_m,
\]
the restriction of $\mat{A}$ to the subspace, expressed in the orthonormal basis. Its
entries are $(\mat{T}_m)_{ij} = \inner{\mat{A}\vq_j}{\vq_i}$---exactly the
orthogonalization coefficients the Arnoldi recurrence already computed and stored
(\cref{ch:orthog}). So $\mat{T}_m$ \emph{is the Hessenberg matrix} $\bar{\mat{H}}_m$
(its leading $m\times m$ block), produced as a free by-product of building the basis. For
symmetric $\mat{A}$ it is the \emph{tridiagonal} matrix of Lanczos. Either way, $\mat{T}_m$
is small---$m$ is a few dozen, not a few million---and dense, so we hand it to the QR
algorithm and get \emph{all} of its eigenpairs cheaply.

\begin{figure}[t]
\centering
\begin{tikzpicture}[font=\footnotesize, scale=1.0]
  % ===== big operator A =====
  \draw[thick, fill=simBgGray] (-0.2,-1.6) rectangle (2.0,1.6);
  \node[font=\small\bfseries, text=simInk] at (0.9,1.95) {$\mat{A}$};
  \node[font=\scriptsize, text=simGray] at (0.9,0.2) {$N\times N$};
  \node[font=\scriptsize, text=simGray] at (0.9,-0.2) {sparse, huge};
  % sparse dots
  \foreach \p in {(0.1,1.2),(0.6,0.9),(1.3,1.4),(1.7,0.5),(0.3,-0.5),(1.0,-1.1),(1.6,-0.8),(0.5,-1.3),(1.4,-0.1)}{
    \fill[simBlue!55] \p circle (1.2pt);}
  % projection arrow
  \draw[flow, simViolet, very thick] (2.3,0) -- (3.6,0)
    node[midway, above, font=\scriptsize, text=simViolet, align=center]
      {$\mat{V}_m^{\!\top}(\cdot)\mat{V}_m$\\project onto $\mathcal{K}_m$};
  % ===== small projected T_m =====
  \draw[thick, fill=simBgBlue] (3.9,-0.7) rectangle (5.3,0.7);
  \node[font=\small\bfseries, text=simBlue] at (4.6,1.0) {$\mat{T}_m$};
  \node[font=\scriptsize, text=simGray] at (4.6,-1.05) {$m\times m$, dense};
  % tridiagonal hint
  \foreach \r in {0,1,2}{\foreach \c in {0,1,2}{
    \pgfmathtruncatemacro\d{abs(\r-\c)}
    \ifnum\d<2 \fill[simBlue!60] (4.0+\c*0.42,0.5-\r*0.42) circle (1.6pt)\fi;}}
  % QR arrow
  \draw[flow, simRust, very thick] (5.6,0) -- (6.9,0)
    node[midway, above, font=\scriptsize, text=simRust, align=center]
      {dense QR\\(cheap)};
  % ===== Ritz pairs =====
  \node[product, align=left, text width=30mm] at (8.5,0)
    {Ritz values\\ $\theta_1,\dots,\theta_m$\\[2pt] Ritz vectors\\ $\mat{V}_m\vs_i$};
  \node[font=\scriptsize\itshape, text=simGreen, align=center, text width=42mm]
    at (5.0,-2.25)
    {the extreme $\theta_i$ approximate $\mat{A}$'s extreme eigenvalues---first and best};
\end{tikzpicture}
\caption[Rayleigh--Ritz projection]{The Rayleigh--Ritz idea in one picture. The huge
sparse operator $\mat{A}$ (left) is projected onto the small Krylov subspace
$\mathcal{K}_m$ spanned by the orthonormal basis $\mat{V}_m$, producing the tiny
\emph{projected operator} $\mat{T}_m = \mat{V}_m^{\!\top}\mat{A}\mat{V}_m$
(centre)---which is exactly the Hessenberg/tridiagonal matrix the Arnoldi/Lanczos
recurrence already filled (\cref{ch:orthog}). Because $\mat{T}_m$ is small and dense, the
classical dense QR algorithm extracts \emph{all} its eigenpairs $(\theta_i, \vs_i)$
cheaply. The \emph{Ritz values} $\theta_i$ and \emph{Ritz vectors} $\mat{V}_m\vs_i$
(right) are the subspace's best approximations to $\mat{A}$'s eigenpairs---and crucially,
the \emph{extreme} ones converge first and most accurately. We never diagonalize
$\mat{A}$; we diagonalize its shadow on a cleverly chosen small subspace.}
\label{fig:rayleighritz}
\end{figure}

\subsection{Ritz values and Ritz vectors}

Let $(\theta_i, \vs_i)$ be an eigenpair of the small matrix $\mat{T}_m$, so
$\mat{T}_m\vs_i = \theta_i\vs_i$ with $\vs_i \in \Real^m$. The scalar $\theta_i$ is a
\emph{Ritz value}, and the vector $\mat{V}_m\vs_i \in \Real^N$---the small eigenvector
$\vs_i$ lifted back into the full space through the basis---is the corresponding
\emph{Ritz vector}. \Cref{fig:rayleighritz} traces the whole pipeline.

\begin{definition}[Ritz pairs]
\label{def:ritz}
Given an orthonormal basis $\mat{V}_m$ of a subspace and the projected operator
$\mat{T}_m = \mat{V}_m^{\!\top}\mat{A}\mat{V}_m$, let $\mat{T}_m\vs_i = \theta_i\vs_i$.
The pair $(\theta_i,\, \mat{V}_m\vs_i)$ is a \emph{Ritz pair}: $\theta_i$ is the
\emph{Ritz value} and $\vx_i := \mat{V}_m\vs_i$ the \emph{Ritz vector}. Ritz pairs are
the subspace's approximations to the eigenpairs of $\mat{A}$; in exact arithmetic, when
$\mathcal{K}_m$ contains a true eigenvector, the corresponding Ritz pair is that
eigenpair exactly.
\end{definition}

The claim that earns the construction its keep is this: the Ritz values approximate the
eigenvalues of $\mat{A}$, and they approximate the \emph{extremal} ones---the largest
and smallest---first and best. Why the extremes? The reason is the same effect that made
the power basis collapse toward the dominant eigenvector (\cref{ch:orthog}), now turned
to advantage. The Krylov subspace $\mathcal{K}_m(\mat{A}, \vv) =
\operatorname{span}\{\vv, \mat{A}\vv, \dots, \mat{A}^{m-1}\vv\}$ is exactly the set of
vectors $p(\mat{A})\vv$ for polynomials $p$ of degree $< m$. Expand the starting vector
in $\mat{A}$'s eigenbasis, $\vv = \sum_k c_k\vu_k$; then $p(\mat{A})\vv = \sum_k
c_k\,p(\lambda_k)\vu_k$. The subspace can therefore reach any combination of eigenvectors
that some degree-$(m{-}1)$ polynomial, evaluated at the eigenvalues, can produce. And a
polynomial of modest degree can take a large value at one extreme eigenvalue while
staying small at all the interior ones---the interior eigenvalues are crowded together,
the extremes stick out. So the subspace captures the extreme eigenvectors with a
low-degree polynomial (small $m$), while the interior ones need a high-degree polynomial
that separates neighbours (large $m$). This is the polynomial-approximation view, and it
is why the extremes come first.

\begin{keyidea}
The Krylov subspace $\mathcal{K}_m(\mat{A},\vv)$ is the span of $\{p(\mat{A})\vv : \deg
p < m\}$. A low-degree polynomial can pick out an \emph{extreme} eigenvalue---one that
sticks out from the rest---while a high-degree one is needed to resolve \emph{interior}
eigenvalues that are crowded together. So Rayleigh--Ritz on a small Krylov subspace
captures the extremal eigenpairs first and most accurately, and reaches the interior
ones only slowly. This is precisely why the eigenmode driver uses the \emph{shift-invert}
spectral transform of \cref{ch:eigenvalues}: it remaps the wanted interior eigenvalues
to the extremes, where Krylov projection is fast (\cref{sec:shiftinvert}).
\end{keyidea}

There is a cleaner way to see ``best approximation'' that the polynomial picture only
gestures at. Among all scalars $\theta$ and all unit vectors $\vx$ drawn from the
subspace, the Ritz values are the stationary values of the \emph{Rayleigh quotient}
$\rho(\vx) = \inner{\mat{A}\vx}{\vx} / \inner{\vx}{\vx}$ restricted to $\mathcal{K}_m$.
For symmetric $\mat{A}$ the Rayleigh quotient's global maximum and minimum over \emph{all
of $\Real^N$} are the largest and smallest eigenvalues (the Courant--Fischer min-max
theorem); restricting the search to the subspace $\mathcal{K}_m$ gives the best bound the
subspace can offer, and the extreme Ritz values are exactly that best bound. The largest
Ritz value is a \emph{lower} bound on the largest eigenvalue, the smallest Ritz value an
\emph{upper} bound on the smallest---the Ritz values bracket the spectrum from inside and
tighten as the subspace grows.

\subsection{The cheap convergence test}

A Ritz pair is an \emph{approximate} eigenpair. We need to know how good it is---and the
beautiful fact, the one that makes the whole method practical, is that we can measure its
quality \emph{without ever forming the Ritz vector or applying $\mat{A}$ to it}. The cost
of the convergence test is essentially zero.

The natural quality measure is the eigen-residual: if $(\theta_i, \vx_i)$ with
$\vx_i = \mat{V}_m\vs_i$ were a true eigenpair, then $\mat{A}\vx_i - \theta_i\vx_i$ would
vanish. So the residual norm $\norm{\mat{A}\vx_i - \theta_i\vx_i}$ measures how far off it
is. Forming this directly would cost a matrix-vector product with $\mat{A}$ on the full
vector $\vx_i$---not free. But the Arnoldi relation $\mat{A}\mat{V}_m = \mat{V}_{m+1}
\bar{\mat{H}}_m$ of \cref{ch:orthog} collapses it to something we can read off the small
data we already have. Substitute $\vx_i = \mat{V}_m\vs_i$ and use the relation:
\[
  \mat{A}\vx_i - \theta_i\vx_i
    = \mat{A}\mat{V}_m\vs_i - \theta_i\mat{V}_m\vs_i
    = \mat{V}_{m+1}\bar{\mat{H}}_m\vs_i - \theta_i\mat{V}_m\vs_i.
\]
Split $\bar{\mat{H}}_m$ into its top $m\times m$ block $\mat{T}_m$ and its single
sub-diagonal spill-over row, whose only nonzero is $h_{m+1,m}$ in the last column. Since
$\mat{T}_m\vs_i = \theta_i\vs_i$, the $\mat{T}_m$ part cancels the $\theta_i\mat{V}_m\vs_i$
term exactly, leaving only the spill-over:
\[
  \mat{A}\vx_i - \theta_i\vx_i \;=\; h_{m+1,m}\,(\ve_m^{\!\top}\vs_i)\,\vq_{m+1},
\]
where $\ve_m^{\!\top}\vs_i$ is the \emph{last component} of the small eigenvector $\vs_i$
and $\vq_{m+1}$ is the next (unit) basis vector. Taking norms, and using
$\norm{\vq_{m+1}} = 1$:
\[
  \boxed{\;\norm{\mat{A}\vx_i - \theta_i\vx_i}
    \;=\; \abs{h_{m+1,m}}\;\abs{\ve_m^{\!\top}\vs_i}\;}
\]
The residual norm of a full-size Ritz pair is the product of \emph{two scalars}: the
sub-diagonal entry $h_{m+1,m}$ (a single number, already computed when the basis was
extended) and the last component of the tiny eigenvector $\vs_i$ (read directly from the
QR solve of $\mat{T}_m$). No full-vector operation, no application of $\mat{A}$. The test
costs nothing.

\begin{figure}[t]
\centering
\begin{tikzpicture}[font=\footnotesize, scale=1.0]
  % H-bar block: T_m on top, spill-over row at bottom
  \draw[thick, fill=simBgBlue] (0,0) rectangle (2.4,2.4);
  \node[font=\scriptsize, text=simBlue] at (1.2,1.5) {$\mat{T}_m$};
  \node[font=\scriptsize, text=simGray] at (1.2,2.05) {(top $m\times m$ block)};
  % spill-over row
  \draw[thick, fill=simBgRust] (0,-0.5) rectangle (2.4,0.0);
  \node[font=\scriptsize, text=simRust] at (1.95,-0.25) {$h_{m+1,m}$};
  \node[font=\scriptsize, text=simGray] at (0.7,-0.25) {$0\;\cdots\;0$};
  \node[font=\small, text=simInk] at (1.2,-0.95) {$\bar{\mat{H}}_m$};
  % eigenvector s_i
  \draw[thick, fill=simBgGreen] (3.1,0) rectangle (3.55,2.4);
  \node[font=\scriptsize, text=simGreen] at (3.32,1.5) {$\vs_i$};
  \draw[thick, fill=simGreen!40] (3.1,0) rectangle (3.55,0.4);
  \node[font=\scriptsize, text=simGreen, anchor=west] at (3.7,0.2)
    {last component $\ve_m^{\!\top}\vs_i$};
  % product = residual
  \node[opbox, align=center, anchor=west] at (5.6,1.7)
    {$\norm{\mat{A}\vx_i - \theta_i\vx_i}$\\$=\;\abs{h_{m+1,m}}\,\abs{\ve_m^{\!\top}\vs_i}$};
  \node[font=\scriptsize\itshape, text=simGray, align=left, anchor=west] at (5.6,0.4)
    {a product of \emph{two scalars}:\\both already computed.\\no full-vector work,\\
     no apply of $\mat{A}$.};
  \draw[refedge, simRust] (2.4,-0.25) .. controls (4.0,-0.25) and (5.0,1.0) .. (5.55,1.55);
  \draw[refedge, simGreen] (3.55,0.2) .. controls (4.6,0.2) and (5.0,1.2) .. (5.55,1.45);
\end{tikzpicture}
\caption[The cheap Ritz residual test]{Why the convergence test is free. The eigen-residual
of a Ritz pair $(\theta_i, \mat{V}_m\vs_i)$ would seem to need a full matrix-vector product
with $\mat{A}$, but the Arnoldi relation collapses it to a product of two already-computed
scalars. The Hessenberg matrix $\bar{\mat{H}}_m$ splits into its top block $\mat{T}_m$
(blue, whose eigenpair is $(\theta_i,\vs_i)$) and a single spill-over row whose only
nonzero is the sub-diagonal entry $h_{m+1,m}$ (rust). Because $\mat{T}_m\vs_i =
\theta_i\vs_i$ cancels the projected part exactly, the residual reduces to $h_{m+1,m}$
times the \emph{last component} $\ve_m^{\!\top}\vs_i$ of the small eigenvector (green)---two
numbers we already hold. The test $\abs{h_{m+1,m}}\abs{\ve_m^{\!\top}\vs_i} < \texttt{tol}$
is what the eigensolver checks every step to decide which Ritz pairs have converged.}
\label{fig:ritzresidual}
\end{figure}

\begin{notationbox}
A Ritz pair has \emph{converged} when its cheap residual falls below the requested
tolerance:
\[
  \abs{h_{m+1,m}}\;\abs{\ve_m^{\!\top}\vs_i} \;<\; \texttt{tol}\cdot\abs{\theta_i}
  \quad\text{(relative)}.
\]
Two readings of this formula are worth keeping. First, a Ritz pair converges when the
\emph{last component} $\ve_m^{\!\top}\vs_i$ of its small eigenvector is tiny---that is,
when the eigenvector lives almost entirely in the earlier basis vectors and barely
touches the newest one. Second, the shared factor $h_{m+1,m}$ multiplies \emph{every}
pair's residual: when the sub-diagonal $h_{m+1,m}$ shrinks toward zero, the basis is
approaching an \emph{invariant subspace} and \emph{all} the Ritz pairs converge at once.
A small $h_{m+1,m}$ is the eigensolver's signal that the subspace has locked onto a true
invariant subspace of $\mat{A}$.
\end{notationbox}

\subsection{A worked Rayleigh--Ritz}

The whole construction fits on one page of arithmetic. We run Rayleigh--Ritz on a small
symmetric matrix, watch the Ritz values bracket the spectrum, and apply the cheap test.

\begin{workedexample}{Rayleigh--Ritz on a $4\times4$ symmetric matrix}{rrritz}
Take the symmetric matrix and unit starting vector
\[
  \mat{A} = \begin{psmallmatrix}
    2 & 1 & 0 & 0\\ 1 & 2 & 1 & 0\\ 0 & 1 & 2 & 1\\ 0 & 0 & 1 & 2
  \end{psmallmatrix},
  \qquad
  \vq_1 = \begin{psmallmatrix} 1\\0\\0\\0 \end{psmallmatrix}.
\]
The true eigenvalues of $\mat{A}$ are $2 + 2\cos(k\pi/5)$ for $k=1,2,3,4$, namely
\[
  \lambda \;\approx\; \{\,0.382,\; 1.382,\; 2.618,\; 3.618\,\},
\]
with extremes $0.382$ and $3.618$. We will build a $3$-dimensional Krylov basis by
Lanczos, form $\mat{T}_3$, and read its Ritz values---then check how well they bracket the
extremes.

\smallskip
\emph{Build the basis (Lanczos, \cref{ch:orthog}).} Step~1:
$\mat{A}\vq_1 = (2,1,0,0)^{\!\top}$, so $\alpha_1 = \inner{\mat{A}\vq_1}{\vq_1} = 2$ and the
residual is $\mat{A}\vq_1 - 2\vq_1 = (0,1,0,0)^{\!\top}$, giving $\beta_1 = 1$ and
$\vq_2 = (0,1,0,0)^{\!\top}$. Step~2: $\mat{A}\vq_2 = (1,2,1,0)^{\!\top}$, so $\alpha_2 = 2$
and the three-term residual is $\mat{A}\vq_2 - 2\vq_2 - 1\cdot\vq_1 = (0,0,1,0)^{\!\top}$,
giving $\beta_2 = 1$ and $\vq_3 = (0,0,1,0)^{\!\top}$. The projected operator is the
tridiagonal
\[
  \mat{T}_3 \;=\; \begin{psmallmatrix}
    \alpha_1 & \beta_1 & 0\\ \beta_1 & \alpha_2 & \beta_2\\ 0 & \beta_2 & \alpha_3
  \end{psmallmatrix}
  \;=\; \begin{psmallmatrix} 2 & 1 & 0\\ 1 & 2 & 1\\ 0 & 1 & 2 \end{psmallmatrix},
\]
where $\alpha_3 = \inner{\mat{A}\vq_3}{\vq_3} = 2$ completes the diagonal. (We also pick up
$\beta_3 = \norm{\mat{A}\vq_3 - 2\vq_3 - \vq_2}$; one computes $\mat{A}\vq_3 = (0,1,2,1)^{\!\top}$
and the residual $(0,0,0,1)^{\!\top}$, so $h_{4,3} = \beta_3 = 1$---the spill-over entry we
need for the residual test.)

\smallskip
\emph{Compute the Ritz values.} The eigenvalues of the $3\times3$ $\mat{T}_3$ are
$2$ and $2 \pm \sqrt2$, i.e.
\[
  \theta \;=\; \{\,0.586,\; 2.000,\; 3.414\,\}.
\]
Compare to the true spectrum. The extreme Ritz values $0.586$ and $3.414$ bracket the
true extremes $0.382$ and $3.618$ \emph{from inside}: $0.586 > 0.382$ and $3.414 < 3.618$,
exactly as the min-max bound promises (the smallest Ritz value over-estimates the smallest
eigenvalue; the largest under-estimates the largest). With only a $3$-dimensional subspace
of a $4$-dimensional space, the extremes are already approximated to within about $0.2$,
while the interior pair is represented by the single middling Ritz value $2.000$.

\smallskip
\emph{Apply the cheap test.} Take the largest Ritz value $\theta = 3.414$ with eigenvector
$\vs = \tfrac12(1, \sqrt2, 1)^{\!\top}$ of $\mat{T}_3$ (the standard eigenvector of
$\mathrm{tridiag}(1;2;1)$). Its last component is $\ve_3^{\!\top}\vs = \tfrac12$. The
residual of the corresponding full Ritz pair is then, with no further matrix work,
\[
  \norm{\mat{A}\vx - \theta\vx} \;=\; \abs{h_{4,3}}\,\abs{\ve_3^{\!\top}\vs}
    \;=\; 1 \cdot \tfrac12 \;=\; 0.5.
\]
This is not yet converged---a residual of $0.5$ is large---which is consistent with the
Ritz value still being off by $0.2$ from the true eigenvalue. One more Lanczos step (here
filling the subspace to all of $\Real^4$) would drive $h_{m+1,m}$ to zero and the residual
with it: at $m = 4 = N$ the subspace is invariant, $h_{5,4} = 0$, and the Ritz values
\emph{equal} the eigenvalues exactly. The cheap test tracked the convergence the whole way
without one extra application of $\mat{A}$.
\end{workedexample}

\Cref{fig:ritzconverge} shows the same story for a larger matrix, where the subspace is
genuinely smaller than the space: the extreme Ritz values march toward the true extremes as
the basis grows, fast at the ends and slow in the middle.

\begin{figure}[t]
\centering
\begin{tikzpicture}
\begin{axis}[
  width=12cm, height=7cm,
  xlabel={basis size $m$},
  ylabel={Ritz values $\theta_i$ \ (vs.\ true eigenvalues, dashed)},
  xmin=1, xmax=12, ymin=-0.3, ymax=8.3,
  grid=both, grid style={simGray!18},
  legend pos=south east, legend cell align=left,
  legend style={font=\scriptsize, fill=white, draw=simGray!40},
  label style={font=\footnotesize}, tick label style={font=\scriptsize},
]
% true extreme eigenvalues (dashed reference lines)
\addplot[simGray, densely dashed, thick, domain=1:12, samples=2] {7.7};
\addplot[simGray, densely dashed, thick, domain=1:12, samples=2] {0.15};
\node[font=\scriptsize, text=simGray, anchor=west] at (axis cs:9.2,7.95) {$\lambda_{\max}$};
\node[font=\scriptsize, text=simGray, anchor=west] at (axis cs:9.2,0.5) {$\lambda_{\min}$};
% largest Ritz value: converges UP to lambda_max quickly
\addplot[simRust, very thick, mark=*, mark size=1.3pt] coordinates {
  (1,4.0)(2,5.6)(3,6.6)(4,7.15)(5,7.45)(6,7.6)(7,7.66)(8,7.69)(9,7.70)(10,7.70)(11,7.70)(12,7.70)};
\addlegendentry{largest Ritz value $\theta_{\max}$}
% smallest Ritz value: converges DOWN to lambda_min quickly
\addplot[simBlue, very thick, mark=square*, mark size=1.3pt] coordinates {
  (1,4.0)(2,2.3)(3,1.3)(4,0.75)(5,0.45)(6,0.30)(7,0.22)(8,0.18)(9,0.16)(10,0.15)(11,0.15)(12,0.15)};
\addlegendentry{smallest Ritz value $\theta_{\min}$}
% an interior Ritz value: converges slowly, wobbles
\addplot[simGreen, very thick, mark=triangle*, mark size=1.5pt] coordinates {
  (3,3.1)(4,3.6)(5,3.9)(6,4.0)(7,3.95)(8,3.92)(9,3.9)(10,3.9)(11,3.9)(12,3.9)};
\addlegendentry{an interior Ritz value (slow)}
\addplot[simGray, densely dotted, thick, domain=1:12, samples=2] {3.9};
\end{axis}
\end{tikzpicture}
\caption[Ritz values converging to the extremes]{How the Ritz values approach the true
eigenvalues as the Krylov basis grows---a schematic of the universal behaviour. The
\emph{extreme} eigenvalues are captured fast: the largest Ritz value (rust) climbs to
$\lambda_{\max}$ and the smallest (blue) descends to $\lambda_{\min}$ within a handful of
steps, each approaching its target monotonically from inside the spectrum (the min-max
bracketing of \cref{sec:rayleighritz}). An \emph{interior} eigenvalue (green) is resolved
only slowly and non-monotonically, because the subspace must build up a high-degree
polynomial to separate it from its crowded neighbours. This asymmetry---extremes fast,
interior slow---is the whole reason the eigenmode driver applies the shift-invert
transform of \cref{ch:eigenvalues} before iterating: it moves the wanted interior modes to
the extremes, where this convergence is fast.}
\label{fig:ritzconverge}
\end{figure}

\section{The restart problem}
\label{sec:restart}

Rayleigh--Ritz says: build the basis, project, read the Ritz values. The hidden assumption
is that we can afford to let the basis grow until the wanted Ritz values converge. For the
extremes on a well-behaved operator that may take only a few dozen steps---but for a
clustered spectrum, or many wanted modes, the basis $\mat{V}_m$ can grow to hundreds or
thousands of columns before convergence. And a growing basis is expensive in two distinct
ways, plus a third that is worse than expensive.

\begin{figure}[t]
\centering
\begin{tikzpicture}
\begin{axis}[
  width=12cm, height=6.3cm,
  xlabel={basis size $m$},
  ylabel={relative cost / loss},
  ymode=log, ymin=1e-15, ymax=3e3,
  xmin=0, xmax=200,
  grid=both, grid style={simGray!18},
  legend pos=outer north east, legend cell align=left,
  legend style={font=\scriptsize, fill=white, draw=simGray!40},
  label style={font=\footnotesize}, tick label style={font=\scriptsize},
]
% storage: linear in m
\addplot[simBlue, very thick, domain=1:200, samples=40] {x};
\addlegendentry{storage $\sim m$ (vectors held)}
% orthogonalization cost per step: linear in m; cumulative quadratic. show per-step ~ m
\addplot[simRust, very thick, domain=1:200, samples=40] {0.5*x};
\addlegendentry{orthog.\ cost/step $\sim m$}
% loss of orthogonality (Lanczos): grows, eventually O(1)
\addplot[simViolet, very thick, domain=1:200, samples=60] {2e-15*exp(0.13*x)};
\addlegendentry{Lanczos loss-of-orthog.\ $\eta_m$}
% the O(1) danger line
\addplot[simGray, densely dashed, thick, domain=0:200, samples=2] {1};
\node[font=\scriptsize, text=simGray, anchor=west] at (axis cs:2,2.3) {orthogonality lost};
% the restart cap
\draw[simGreen, very thick, densely dashed] (axis cs:40,1e-15) -- (axis cs:40,3e3);
\node[font=\scriptsize\bfseries, text=simGreen, rotate=90, anchor=south]
  at (axis cs:40,3e-7) {restart here ($m=k_{\max}$)};
\end{axis}
\end{tikzpicture}
\caption[The restart problem]{Why the basis cannot be allowed to grow without bound. As the
basis size $m$ increases, three costs climb. \emph{Storage} (blue) grows linearly---every
basis vector is a full field, and a thousand of them may not fit in memory. \emph{Orthogonalization
cost} (rust): each new Arnoldi column must be orthogonalized against all $m$ previous ones, so
the per-step cost grows linearly in $m$ and the cumulative cost \emph{quadratically}. Worst of
all, for Lanczos the \emph{loss of orthogonality} (violet, \cref{ch:orthog}) grows until the
basis is no longer orthogonal at all (dashed line at $\eta_m \sim 1$), at which point the
projected $\mat{T}_m$ reports \emph{ghost eigenvalues}---spurious copies of converged
eigenvalues that are artifacts, not real. The cure is to cap the basis at some $k_{\max}$
(green) and \emph{restart}---but a naive restart throws away the converged Ritz vectors. The
art of \cref{sec:thickrestart} is to restart \emph{without} discarding the good directions.}
\label{fig:restartproblem}
\end{figure}

\Cref{fig:restartproblem} tabulates the three pressures. The first is \emph{storage}: each
basis vector $\vq_j$ is a full field of $N$ entries, and $m$ of them occupy $mN$ numbers. At a
million unknowns, a thousand basis vectors is a billion numbers---gigabytes, possibly more than
fits in memory. The second is \emph{orthogonalization cost}: Arnoldi orthogonalizes each new
column against all $m$ previous columns (\cref{ch:orthog}), so step $m$ costs $O(m)$ inner
products and the whole build costs $O(m^2)$. The basis that took $O(m)$ to store takes $O(m^2)$
to build. The third pressure is the dangerous one, and it is specific to Lanczos: \emph{loss of
orthogonality}. The three-term recurrence trusts that its projections against the old basis
vectors are zero and never re-checks them; in finite precision they are not zero, the error
compounds, and the basis silently develops near-duplicate directions
(\cref{ch:orthog}). The projected $\mat{T}_m$ then reports \emph{ghost eigenvalues}---the same
eigenvalue several times over, spurious copies that a naive solver mistakes for a genuine
high-multiplicity mode.

So the basis must be capped. Pick a maximum size $k_{\max}$, build up to it, and then---do what,
exactly? This is the restart problem, and the naive answer is bad.

\begin{pitfall}
The naive restart \emph{discards the basis and starts over} from a single new vector---chosen,
say, as the best converged Ritz vector so far, in the hope of steering the next subspace toward
the wanted modes. This throws away almost everything you computed. The basis $\mat{V}_{k_{\max}}$
encoded \emph{several} nearly-converged Ritz directions, each the fruit of dozens of operator
applications and inner solves; collapsing it to one vector keeps at most one of them and forces
the others to be rediscovered from scratch. Worse, the restart is repeated every cycle, so the
waste compounds: the method spends most of its operator-applies re-converging directions it had
already converged and then threw away. A naive-restart eigensolver can be an order of magnitude
slower than one that restarts intelligently. The fix is to keep \emph{all} the good directions,
not just one---which is exactly what thick-restart does.
\end{pitfall}

The tension is now sharp. We must keep the basis small (storage, cost, orthogonality all demand
it), but we must not throw away the converged and nearly-converged Ritz directions (they are the
expensive part of the answer). A restart that resolves both---small basis, good directions
retained---is what we need, and it exists.

\section{Thick-restart and Krylov--Schur}
\label{sec:thickrestart}

The resolution is due to Stewart, building on the earlier thick-restart Lanczos of Wu and
Simon: restart from a \emph{small} basis that still spans all the Ritz directions worth keeping.
Instead of collapsing $\mat{V}_{k_{\max}}$ to a single vector, we compress it to a basis of
size $k$ (a few more than the number of wanted modes) that retains the $k$ best Ritz directions
and discards the other $k_{\max} - k$. The trick is to do this compression in a way that
\emph{preserves the Krylov structure}---so that the next expansion phase is a legitimate Arnoldi
continuation, not a fresh start.

\subsection{The partial Schur form}

The key is a change of representation. The Arnoldi relation $\mat{A}\mat{V}_m =
\mat{V}_{m+1}\bar{\mat{H}}_m$ stores the projected operator as a Hessenberg matrix $\mat{T}_m$.
Krylov--Schur stores it instead as a \emph{partial Schur form}: it computes the Schur
decomposition $\mat{T}_m = \mat{S}\,\mat{R}\,\mat{S}^{\!\top}$ (with $\mat{R}$ quasi-triangular,
$\mat{S}$ orthogonal) and \emph{reorders} it so that the wanted Ritz values---the converged and
promising ones---sit in the leading $k \times k$ block. Transforming the basis by $\mat{S}$,
$\tilde{\mat{V}}_m = \mat{V}_m\mat{S}$, rotates the orthonormal basis so that its first $k$
columns span exactly the wanted Ritz directions and the projected operator on them is the clean
triangular block $\mat{R}_{1:k,1:k}$.

The structural fact that makes this a \emph{restart} rather than just a reordering is that the
Arnoldi relation survives the transformation in a slightly generalized form. After reordering,
\[
  \mat{A}\,\tilde{\mat{V}}_k
    \;=\; \tilde{\mat{V}}_k\,\mat{R}_{1:k,1:k}
       \;+\; \vq_{m+1}\,\vb_k^{\!\top},
\]
where $\tilde{\mat{V}}_k$ is the first $k$ transformed columns, $\mat{R}_{1:k,1:k}$ the leading
triangular block, and $\vb_k^{\!\top}$ a short row vector (the transformed spill-over). This is a
\emph{Krylov--Schur relation}: like the Arnoldi relation it expresses $\mat{A}$'s action on a
$k$-column basis in terms of that basis plus a single extra direction $\vq_{m+1}$---which means
we can \emph{continue} the Arnoldi expansion from here, appending new columns by the ordinary
recurrence, exactly as if $k$ steps had just been taken. The compressed basis is a valid
launching point for more Arnoldi steps. That is the whole game: compress to size $k$, keeping the
good directions, in a form from which expansion can resume.

\begin{figure}[t]
\centering
\begin{tikzpicture}[font=\footnotesize, scale=0.92]
  % ===== before: full basis of size k_max =====
  \begin{scope}
    \node[font=\small\bfseries, text=simInk] at (1.4,3.1) {before restart: size $k_{\max}$};
    % a row of basis columns, some "good" (green), some "junk" (gray)
    \foreach \i/\col/\lab in {0/simGreen/g,1/simGreen/g,2/simGreen/g,3/simGray/j,4/simGray/j,5/simGray/j,6/simGray/j}{
      \draw[thick, fill=\col!25, draw=\col] (\i*0.42,0) rectangle (\i*0.42+0.36,2.4);
    }
    \node[font=\scriptsize, text=simGreen] at (0.55,-0.4) {$k$ good Ritz dirs};
    \node[font=\scriptsize, text=simGray] at (2.3,-0.4) {$k_{\max}-k$ junk};
    \draw[decorate, decoration={brace, amplitude=4pt}, simGreen] (-0.05,2.55) -- (1.3,2.55);
    \draw[decorate, decoration={brace, amplitude=4pt}, simGray] (1.32,2.55) -- (2.94,2.55);
  \end{scope}
  % ===== arrow: thick-restart compression =====
  \draw[flow, simViolet, very thick] (3.3,1.2) -- (4.7,1.2)
    node[midway, above, font=\scriptsize, text=simViolet, align=center]
      {Schur reorder\\+ truncate};
  \node[font=\scriptsize\itshape, text=simRust, align=center] at (4.0,0.3)
    {discard junk,\\keep good};
  % ===== after: compressed basis of size k + 1 spill =====
  \begin{scope}[xshift=5.2cm]
    \node[font=\small\bfseries, text=simInk] at (1.0,3.1) {after restart: size $k$};
    \foreach \i in {0,1,2}{
      \draw[thick, fill=simGreen!25, draw=simGreen] (\i*0.42,0) rectangle (\i*0.42+0.36,2.4);}
    % the carried spill-over vector q_{m+1}
    \draw[thick, fill=simBgRust, draw=simRust] (1.55,0) rectangle (1.91,2.4);
    \node[font=\scriptsize, text=simRust, rotate=90] at (1.73,1.2) {$\vq_{m+1}$};
    \node[font=\scriptsize, text=simGreen] at (0.55,-0.4) {$k$ kept dirs};
    \node[font=\scriptsize, text=simRust] at (1.73,-0.75) {carry};
    % expansion resumes
    \draw[flow, simBlue, thick] (2.3,1.2) -- (3.3,1.2)
      node[midway, above, font=\scriptsize, text=simBlue]{expand};
    \foreach \i in {0,1}{
      \draw[thick, fill=simBgBlue, draw=simBlue, densely dashed] (3.5+\i*0.42,0) rectangle (3.5+\i*0.42+0.36,2.4);}
    \node[font=\scriptsize, text=simBlue] at (3.9,-0.4) {new Arnoldi cols};
  \end{scope}
\end{tikzpicture}
\caption[Thick-restart / Krylov--Schur]{One thick-restart cycle. \emph{Left:} the basis has
grown to its cap $k_{\max}$; some columns span converged or promising Ritz directions (green,
the $k$ we want to keep), the rest are exhausted Krylov directions of no further use (grey).
A naive restart would keep only one green column; thick-restart keeps \emph{all $k$}.
\emph{Centre:} the compression reorders the Schur form of the projected operator so the wanted
Ritz values lead, then truncates to the leading $k$ columns---discarding the grey junk while
retaining every good direction. \emph{Right:} the compressed size-$k$ basis, together with the
single carried spill-over vector $\vq_{m+1}$ (rust) from the Krylov--Schur relation, is a valid
Arnoldi launching point: ordinary expansion (blue, dashed) resumes from it, rebuilding toward
$k_{\max}$ before the next compression. The cycle---expand, extract, lock, compress---repeats
until the wanted modes converge.}
\label{fig:thickrestart}
\end{figure}

\subsection{The restart cycle}

With the partial Schur form in hand, the Krylov--Schur eigensolver is a loop of four stages,
shown in \cref{fig:thickrestart} and \cref{fig:ksloop}.

\begin{enumerate}
  \item \textbf{Expand.} From the current basis of size $k$, take Arnoldi (or Lanczos) steps
  until the basis reaches the cap $k_{\max}$, filling the projected operator's new columns. Each
  step is one \code{arnoldi\_step} or \code{lanczos\_step} of \cref{ch:orthog}---and, in the
  eigenmode driver, each contains one inner shift-invert solve (\cref{sec:shiftinvert}).
  \item \textbf{Extract.} Form the projected operator $\mat{T}_{k_{\max}}$, compute its Schur
  form, and read the Ritz values and the cheap residuals $\abs{h_{m+1,m}}\abs{\ve_m^{\!\top}\vs_i}$
  of \cref{sec:rayleighritz}.
  \item \textbf{Lock.} Mark the Ritz pairs whose residual is below tolerance as
  \emph{converged}: they are part of the answer and are frozen, no longer to be touched by the
  iteration. When the locked count reaches the number of requested modes, the loop terminates.
  \item \textbf{Compress (thick-restart).} Reorder the Schur form so the converged and most
  promising Ritz values lead, truncate the basis to the leading $k$ columns, carry the single
  spill-over vector, and return to the Expand stage.
\end{enumerate}

\begin{figure}[t]
\centering
\begin{tikzpicture}[font=\footnotesize, node distance=9mm]
  \node[stage, minimum width=22mm] (expand) {\textbf{Expand}\\[1pt]\scriptsize Arnoldi/Lanczos\\\scriptsize to $k_{\max}$};
  \node[stage, minimum width=22mm, right=18mm of expand] (extract) {\textbf{Extract}\\[1pt]\scriptsize Rayleigh--Ritz\\\scriptsize + cheap residual};
  \node[stage, minimum width=22mm, right=18mm of extract] (lock) {\textbf{Lock}\\[1pt]\scriptsize freeze converged\\\scriptsize Ritz pairs};
  \node[stage, minimum width=22mm, below=14mm of extract] (compress) {\textbf{Compress}\\[1pt]\scriptsize Schur reorder\\\scriptsize truncate to $k$};
  \node[product, right=16mm of lock] (done) {converged\\modes};
  \draw[flow] (expand) -- (extract);
  \draw[flow] (extract) -- (lock);
  \draw[flow] (lock) -- node[above,font=\scriptsize]{all wanted?} (done);
  \draw[flow] (lock.south) |- node[pos=0.25,right,font=\scriptsize]{not yet} (compress.east);
  \draw[flow] (compress) -| (expand.south);
  % inner-solve callout on Expand
  \node[extern, below=10mm of expand, align=center] (inner)
    {\scriptsize each step:\\\scriptsize inner shift-invert\\\scriptsize \code{ksp\_solve} (\cref{ch:eigenvalues})};
  \draw[refedge] (inner.north) -- (expand.south);
\end{tikzpicture}
\caption[The Krylov--Schur restart loop]{The four-stage Krylov--Schur cycle, the algorithm
inside the opaque eigensolver call. \textbf{Expand} the basis by Arnoldi/Lanczos steps up to the
cap $k_{\max}$; \textbf{extract} the Ritz pairs and their cheap residuals by Rayleigh--Ritz;
\textbf{lock} the converged pairs (freezing them out of further iteration), and if all requested
modes are locked, return them as the answer; otherwise \textbf{compress} by thick-restart---reorder
the Schur form to put the good Ritz directions first, truncate to size $k$, and loop back to
expand. The dashed callout records the stack: in the eigenmode driver every Expand step contains
an inner shift-invert linear solve (\cref{sec:shiftinvert}), so the outer loop you see here sits
atop the inner Krylov solver of \cref{ch:multigrid}. This loop is exactly what SLEPc's
\code{EPSKRYLOVSCHUR} runs.}
\label{fig:ksloop}
\end{figure}

This is the constructive driver that the source spec reconstructs from the firm
\code{arnoldi\_step} and \code{lanczos\_step} kernels: an outer thick-restart loop over restart
cycles, wrapping an inner basis-extension loop over basis columns, with a Rayleigh--Ritz
extraction and a lock-and-compress between cycles. In the calculus of \cref{ch:fold} it is two
nested \code{iterate\_while} folds---the outer over restart cycles, the inner over basis
columns---and the body of the inner fold is one Arnoldi/Lanczos step.

\begin{lstlisting}[style=l4style]
-- the Krylov-Schur driver: outer thick-restart fold over an inner basis-extension fold
eigsolve_ks :: LinOp -> EigControl -> EigResult
eigsolve_ks A control =
  let st0 = seed_basis A control                       -- BV <- [v0 / ||v0||]
      loop st = let ext      = extend_basis A st        -- INNER: Arnoldi/Lanczos to k_max
                    (theta, Y) = rayleigh_ritz A ext    -- project + dense QR -> Ritz pairs
                    locked   = lock_converged ext theta Y control  -- cheap residual test
                in thick_restart locked theta Y control  -- reorder Schur, truncate to k
      st_final = iterate_while loop (\st -> not (enough_converged st control)) st0
  in extract_eigpairs A st_final                        -- lift Ritz vecs, un-transform lambda
  where
    -- INNER basis-extension: ncv steps, each ONE arnoldi/lanczos step
    extend_basis A st =
      iterate_while (\s -> append_column s (one_step A s)) (\s -> s.j < control.k_max) st
    one_step A s = if A.hermitian then lanczos_step A s   -- symmetric three-term recurrence
                                  else arnoldi_step A s    -- full Arnoldi
\end{lstlisting}

\begin{workedexample}{One thick-restart step, with small numbers}{thickrestart}
We make the compression concrete on a schematic example. Suppose we have expanded to a basis of
size $k_{\max} = 6$ and want $2$ eigenvalues, restarting to size $k = 3$ (one extra beyond the
wanted count, for headroom). The Rayleigh--Ritz extraction on $\mat{T}_6$ returns six Ritz values,
which we sort by how close they are to the wanted target (here, the largest two), and tag with
their cheap residuals $r_i = \abs{h_{7,6}}\abs{\ve_6^{\!\top}\vs_i}$:
\[
  \begin{array}{c|cccccc}
    \text{Ritz value } \theta_i & 7.91 & 7.40 & 6.2 & 5.1 & 3.8 & 1.2 \\\hline
    \text{residual } r_i & 4\!\times\!10^{-9} & 3\!\times\!10^{-4} & 0.2 & 0.5 & 0.7 & 0.9 \\
    \text{status} & \text{\textbf{locked}} & \text{promising} & \text{keep} & \text{junk} & \text{junk} & \text{junk}
  \end{array}
\]
\emph{Lock.} The first Ritz value $\theta_1 = 7.91$ has residual $4\times10^{-9}$, below the
tolerance $10^{-8}$: it is \emph{converged} and locked---one of the two wanted modes, now frozen.
The second, $\theta_2 = 7.40$ at residual $3\times10^{-4}$, is not yet converged but is the next
wanted mode and converging well: \emph{promising}.

\smallskip
\emph{Choose what to keep.} We keep $k = 3$ directions: the locked $\theta_1$, the promising
$\theta_2$, and the next-best $\theta_3 = 6.2$ for headroom (a converging interior direction worth
continuing). The bottom three Ritz values $\{5.1, 3.8, 1.2\}$---large residuals, far from the
target---are \emph{junk}: exhausted Krylov directions with nothing left to give. They are discarded.

\smallskip
\emph{Compress.} Reorder the Schur form of $\mat{T}_6$ so the kept Ritz values
$\{7.91, 7.40, 6.2\}$ occupy its leading $3\times3$ block, transform the basis
$\tilde{\mat{V}}_6 = \mat{V}_6\mat{S}$, and truncate to the first $3$ columns
$\tilde{\mat{V}}_3$. Together with the carried spill-over vector $\vq_7$, these satisfy the
Krylov--Schur relation $\mat{A}\tilde{\mat{V}}_3 = \tilde{\mat{V}}_3\mat{R}_{1:3} +
\vq_7\vb_3^{\!\top}$, so the next cycle resumes Arnoldi from a \emph{size-$3$} basis that already
spans the converged mode, the next mode, and a promising third---none of which had to be
rediscovered. The basis shrank from $6$ back to $3$, storage and orthogonalization cost reset to
the floor, and \emph{not one} of the good directions was lost. Contrast the naive restart, which
would have kept only $\theta_1$ and forced $\theta_2$ and $\theta_3$ to be rebuilt from a single
vector. That saved rediscovery, multiplied over many cycles, is the entire value of
thick-restart.
\end{workedexample}

\section{Lanczos specifics: the symmetric three-term recurrence}
\label{sec:lanczosspecific}

The eigenmode and boundary-mode analyses pose \emph{symmetric} (or Hermitian) eigenproblems---the
operators $\mat{K}$ and $\mat{M}$ are symmetric positive definite. For these, the Arnoldi basis
extension of the previous sections specializes to the \emph{Lanczos} three-term recurrence of
\cref{ch:orthog}, and the specialization is worth dwelling on because it is both a major economy
and a major hazard.

\subsection{Why symmetric is cheaper}

Recall from \cref{ch:orthog} the symmetry miracle: when $\mat{A} = \mat{A}^{\!\top}$, the projected
Hessenberg matrix $\mat{T}_m = \mat{V}_m^{\!\top}\mat{A}\mat{V}_m$ is itself symmetric, and a
symmetric Hessenberg matrix is \emph{tridiagonal}. So orthogonalizing the new vector
$\mat{A}\vq_j$ against the basis---which for general Arnoldi means projecting against \emph{all}
$j$ previous columns---collapses to projecting against just the \emph{two} most recent,
$\vq_j$ and $\vq_{j-1}$. The projections against $\vq_1, \dots, \vq_{j-2}$ are all exactly zero.
The step is
\[
  \vw \;=\; \mat{A}\vq_j - \alpha_j\vq_j - \beta_{j-1}\vq_{j-1},
  \qquad
  \alpha_j = \inner{\mat{A}\vq_j}{\vq_j},
  \qquad
  \beta_j = \norm{\vw},
  \qquad
  \vq_{j+1} = \frac{\vw}{\beta_j},
\]
storing only $\vq_j, \vq_{j-1}$ and producing the tridiagonal entries $(\alpha_j, \beta_j)$
directly. This is the \code{lanczos\_step} kernel from which the symmetric arm of the eigensolver
is built:

\begin{lstlisting}[style=l4style]
-- the symmetric three-term recurrence: O(1) memory per step, only TWO stored vectors
lanczos_step :: LinOp -> Tensor[$S] -> Tensor[$S] -> Scalar -> (Tensor[$S], Scalar, Scalar)
lanczos_step A q_j q_prev beta_prev =
  let u      = apply A q_j                  -- A q_j
      alpha  = dot u q_j                     -- diagonal coeff alpha_j = <A q_j, q_j>
      w      = axpby (-alpha) q_j 1 u        -- subtract the q_j component
      w'     = axpy (-beta_prev) q_prev w    -- subtract the q_{j-1} component (only TWO!)
      beta   = nrm2 w'                       -- new off-diagonal / normalization beta_j
      q_next = scal (1/beta) w'
  in (q_next, alpha, beta)
\end{lstlisting}

The economy is dramatic. General Arnoldi stores the whole growing basis ($O(m)$ vectors) and
spends $O(m)$ inner products per step; Lanczos stores \emph{two} vectors and spends a constant
number of inner products per step, regardless of $m$. For a symmetric problem the per-step cost
does not grow at all. This is the same short-recurrence economy that gives conjugate gradients its
efficiency (\cref{ch:cg}): symmetry makes a quantity that would otherwise depend on the entire
history depend only on the last two steps.

\subsection{The hazard and the reorthogonalization policy}

The economy is bought on credit, and \cref{ch:orthog} showed the bill: the short recurrence is
short \emph{because} it trusts that the skipped projections are zero, and in finite precision they
are not. Lanczos never looks at the old basis vectors, so it cannot notice---let alone
repair---the orthogonality it loses against them. The orthogonality decays (the violet curve of
\cref{fig:restartproblem}), and the symptom is the \emph{ghost eigenvalue}: the tridiagonal
$\mat{T}_m$ reports a converged eigenvalue several times over, spurious copies of a real eigenvalue
that the lost orthogonality manufactured. A naive Lanczos solver returns a multiplicity-three
eigenvalue that truly has multiplicity one.

The practical eigensolver therefore augments the bare recurrence with a \emph{reorthogonalization
policy}---a decision about when, and against what, to do the extra orthogonalization the short
recurrence omits. \Cref{tab:reorthog} lays out the three standard policies and the cost/robustness
trade each strikes.

\begin{table}[t]
\centering
\caption[Lanczos reorthogonalization policies]{The reorthogonalization policies that buy back the
orthogonality the three-term recurrence loses. \emph{Full} reorthogonalization is the safest and
most expensive---it re-projects every new vector against the entire stored basis, sacrificing the
$O(1)$-memory advantage that motivated Lanczos in the first place. \emph{Selective} and
\emph{partial} reorthogonalization monitor a cheap estimate of the orthogonality loss and
reorthogonalize only when it crosses a threshold, recovering most of the safety at a fraction of
the cost. Which policy the solver runs is a central configuration decision; in a thick-restart
method the restart itself caps the basis size, which bounds the orthogonality loss between restarts
and makes full reorthogonalization affordable.}
\label{tab:reorthog}
\small
\setlength{\tabcolsep}{6pt}
\renewcommand{\arraystretch}{1.25}
\begin{tabular}{@{}llll@{}}
\toprule
Policy & When to reorthogonalize & Against & Cost \\
\midrule
None (bare) & never & --- & cheapest; ghost eigenvalues \\
Full & every step & whole stored basis & most robust; $O(m)$/step, $O(m)$ storage \\
Selective & monitored loss high & converged Ritz vectors only & moderate \\
Partial & monitored loss $>\sqrt{u}$ & recent basis window & moderate; near-full safety \\
\bottomrule
\end{tabular}
\end{table}

\begin{pitfall}
It is tempting to read ``Lanczos is the cheap symmetric method'' and reach for the bare three-term
recurrence whenever the operator is symmetric. Do not, for an eigensolver. Without a
reorthogonalization policy the bare recurrence \emph{will} produce ghost eigenvalues on any
problem run long enough, and the failure is silent: the solver returns confident, wrong
multiplicities with no error flag. The three-term recurrence is the right \emph{kernel}, but it is
not a self-sufficient eigensolver---it must be wrapped in a reorthogonalization policy (and, in
practice, a restart that caps how far orthogonality can decay between restarts). The source spec
records exactly this: a firm Lanczos eigensolver ``must pin a reorthogonalization policy; the
band-3 recurrence alone is not numerically self-sufficient for a converged eigensolve.''
\end{pitfall}

Thick-restart and reorthogonalization are complementary cures for the same disease. The restart
caps the basis at $k_{\max}$, which bounds how far orthogonality can decay before the basis is
compressed and effectively refreshed; reorthogonalization repairs the decay within a single
expansion phase. Together they make symmetric Krylov--Schur both cheap (short recurrence, small
basis) and trustworthy (no ghosts)---the combination the eigenmode driver actually runs.

\section{Putting it together: the shift-invert stack}
\label{sec:shiftinvert}

We have built the eigensolver in the abstract, with $\mat{A}$ standing for ``the operator we
iterate against.'' In the eigenmode driver, $\mat{A}$ is not the original operator. Recall from
\cref{sec:rayleighritz} that Krylov projection captures the \emph{extreme} eigenvalues fast and
the interior ones slowly---but the wanted resonant modes are interior, clustered near a target
frequency, not at the extremes of the spectrum. Running Krylov--Schur on $\mat{K}\vx = \lambda
\mat{M}\vx$ directly would converge agonizingly slowly to exactly the eigenvalues we care about.

The remedy is the \emph{shift-invert} spectral transform of \cref{ch:eigenvalues}. Instead of
iterating against the original pencil, we iterate against
\[
  \mat{A} \;=\; (\mat{K} - \sigma\mat{M})^{-1}\mat{M},
\]
whose eigenvalues are $1/(\lambda - \sigma)$. This transform turns the eigenvalues $\lambda$ near
the shift $\sigma$ into the eigenvalues $1/(\lambda - \sigma)$ of \emph{largest magnitude}---the
extremes. The wanted interior modes become the dominant modes of the transformed operator, and
Krylov--Schur converges to them fast. After the iteration, each converged Ritz value $\theta$ is
un-transformed back to an original-problem eigenvalue $\lambda = \sigma + 1/\theta$.

The crucial structural consequence is that \emph{applying the transformed operator $\mat{A}$
requires an inner linear solve}. There is no matrix $(\mat{K} - \sigma\mat{M})^{-1}$; there is only
the \emph{action} of solving a linear system with $\mat{K} - \sigma\mat{M}$. So one application of
$\mat{A}$ to a vector $\vv$ is two stages: first apply $\mat{M}$ (a matrix-vector product), then
solve $(\mat{K} - \sigma\mat{M})\vy = \mat{M}\vv$ (a full linear solve). That inner solve is the
\code{ksp\_solve} of \cref{ch:krylov}---a Krylov solver in its own right, preconditioned by the
multigrid of \cref{ch:multigrid}, applying the matrix-free operator of \cref{ch:matrixfree}. The
eigensolver's every basis-extension step contains, buried inside its single ``apply $\mat{A}$,'' a
complete inner Krylov solve.

\begin{lstlisting}[style=l4style]
-- one application of the shift-inverted operator: an INNER linear solve inside every Arnoldi step
apply_shift_invert :: EigOp -> Tensor[$S] -> Tensor[$S]
apply_shift_invert op v =
  let w = apply_linop op.operand v     -- stage 1: apply M (or K for no transform)
      y = ksp_solve   op.inv     w     -- stage 2: solve (K - sigma M) y = w  -- INNER Krylov solve!
  in scale_untransform op y            -- per-backend Higham gamma/delta un-scale
\end{lstlisting}

This is the full stack, and it is worth seeing all at once. \Cref{fig:eigenstack} draws it: the
outer Krylov--Schur eigen-iteration owns the basis and the restart; each of its steps applies the
shift-inverted operator; that application contains an inner Krylov linear solve; the inner solve is
accelerated by a preconditioner; and the preconditioner, like the operator, is applied
matrix-free. Five layers, each a chapter of this Part, nested one inside the next.

\begin{figure}[t]
\centering
\begin{tikzpicture}[font=\footnotesize]
  % nested boxes, outer to inner
  \draw[thick, fill=simBgGold, rounded corners=3pt] (-0.2,-0.2) rectangle (10.6,5.4);
  \node[anchor=north west, font=\small\bfseries, text=simGold!60!black] at (0.0,5.3)
    {\,1.\ Krylov--Schur eigen-iteration\quad\scriptsize(this chapter: basis, restart, Ritz)};
  \draw[thick, fill=simBgViolet, rounded corners=3pt] (0.3,0.2) rectangle (10.1,4.5);
  \node[anchor=north west, font=\small\bfseries, text=simViolet] at (0.5,4.4)
    {\,2.\ shift-invert apply $\;\mat{A}=(\mat{K}-\sigma\mat{M})^{-1}\mat{M}$\quad\scriptsize(\cref{ch:eigenvalues})};
  \draw[thick, fill=simBgBlue, rounded corners=3pt] (0.8,0.6) rectangle (9.6,3.6);
  \node[anchor=north west, font=\small\bfseries, text=simBlue] at (1.0,3.5)
    {\,3.\ inner Krylov solve $\;(\mat{K}-\sigma\mat{M})\vy=\vw$\quad\scriptsize\code{ksp\_solve} (\cref{ch:krylov})};
  \draw[thick, fill=simBgGreen, rounded corners=3pt] (1.3,1.0) rectangle (9.1,2.7);
  \node[anchor=north west, font=\small\bfseries, text=simGreen] at (1.5,2.6)
    {\,4.\ preconditioner $\;\mat{P}^{-1}$\quad\scriptsize geometric multigrid (\cref{ch:multigrid})};
  \draw[thick, fill=simBgRust, rounded corners=3pt] (1.8,1.4) rectangle (8.6,2.1);
  \node[anchor=west, font=\small\bfseries, text=simRust] at (2.0,1.75)
    {\,5.\ matrix-free operator apply $\;\mat{K}\vv,\ \mat{M}\vv$\quad\scriptsize(\cref{ch:matrixfree})};
\end{tikzpicture}
\caption[The full eigen-solve stack]{The five nested layers of an eigenmode solve, each a chapter
of this Part. \textbf{(1)} The outer Krylov--Schur eigen-iteration of this chapter owns the basis,
the Rayleigh--Ritz extraction, and the thick-restart---the loop SLEPc calls \code{EPSKRYLOVSCHUR}.
\textbf{(2)} Each basis-extension step applies the shift-inverted operator $\mat{A} = (\mat{K} -
\sigma\mat{M})^{-1}\mat{M}$ of \cref{ch:eigenvalues}, which remaps the wanted interior modes to the
fast-converging extremes. \textbf{(3)} Applying that operator requires an inner Krylov linear solve
$(\mat{K} - \sigma\mat{M})\vy = \vw$---the \code{ksp\_solve} of \cref{ch:krylov}, a solver nested
inside the eigensolver. \textbf{(4)} That inner solve is accelerated by a preconditioner, the
geometric multigrid of \cref{ch:multigrid}. \textbf{(5)} And the operator and preconditioner are
both applied matrix-free (\cref{ch:matrixfree}), never assembled as explicit matrices. The opaque
``eigenvalue solve'' of \cref{ch:eigenvalues} is this entire tower, and every layer is something we
have built.}
\label{fig:eigenstack}
\end{figure}

\begin{keyidea}
The eigenmode solve is a five-layer stack: \emph{Krylov--Schur eigen-iteration} (this chapter)
$\triangleright$ \emph{shift-invert apply} (\cref{ch:eigenvalues}) $\triangleright$ \emph{inner
Krylov solve} (\cref{ch:krylov}) $\triangleright$ \emph{preconditioner} (\cref{ch:multigrid})
$\triangleright$ \emph{matrix-free operator apply} (\cref{ch:matrixfree}). The outer eigen-loop is
what the library owns; the inner linear solve is the \code{ksp\_solve} the rest of this Part built.
Every basis-extension step of the eigensolver contains one full inner Krylov solve, so the cost of
an eigenmode analysis is dominated by the inner solves---which is why the preconditioner of
\cref{ch:multigrid} matters as much to the eigensolver as to any linear solve.
\end{keyidea}

\section{The library boundary, demystified}
\label{sec:libraryboundary}

We can now return to the promise of \cref{ch:eigenvalues}: that the eigenmode driver's
``solve''---the one stage of the whole simulator that was a single call into an opaque library
rather than a loop we wrote---would turn out to be a loop we could have written ourselves. We have
now written it. The four-stage cycle of \cref{fig:ksloop}, with the Lanczos or Arnoldi kernel of
\cref{ch:orthog} at its heart and the shift-invert stack of \cref{fig:eigenstack} inside each step,
\emph{is} the algorithm that SLEPc's \code{EPSSolve} with the \code{EPSKRYLOVSCHUR} type runs, and
that ARPACK runs through its reverse-communication \code{naupd}.

The phrase ``reverse communication'' names the only genuinely unusual thing about the library
boundary, and it is a packaging detail, not an algorithmic one. ARPACK does not call Palace's
operator; it cannot, because it was compiled with no knowledge of Palace's data structures.
Instead, the control inverts: Palace calls into ARPACK's \code{naupd} driver, which runs a little
of the Krylov--Schur loop and then \emph{returns}, setting a flag that says ``I need
$\mat{A}\vv$ for this vector $\vv$---compute it and call me again.'' Palace computes the matrix-vector
product (running the entire shift-invert stack of \cref{sec:shiftinvert} to do so) and calls
\code{naupd} again; the driver resumes where it left off, takes the next step, and returns with the
next request. The eigen-iteration logic lives entirely inside the library; Palace's loop is a
dispatcher that services the library's requests for operator applications.

\begin{figure}[t]
\centering
\begin{tikzpicture}[font=\footnotesize, node distance=10mm]
  % Palace side
  \node[stage, minimum width=30mm, align=center] (palace)
    {\textbf{Palace}\\[1pt]\scriptsize dispatcher loop\\\scriptsize + operator applies};
  % library side
  \node[extern, minimum width=34mm, minimum height=20mm, right=34mm of palace, align=center] (lib)
    {\textbf{SLEPc / ARPACK}\\[2pt]\scriptsize \emph{the Krylov--Schur loop}\\\scriptsize basis $\cdot$ restart $\cdot$ Ritz\\[3pt]
     \scriptsize\itshape (this chapter's\\\scriptsize\itshape algorithm, inside)};
  % the two reverse-communication arrows
  \draw[flow, simViolet] (palace.north east) .. controls +(1.0,0.8) and +(-1.0,0.8) .. (lib.north west)
    node[midway, above, font=\scriptsize, text=simViolet, align=center]{call \code{naupd};\\``take a step''};
  \draw[backflow] (lib.south west) .. controls +(-1.0,-0.8) and +(1.0,-0.8) .. (palace.south east)
    node[midway, below, font=\scriptsize, text=simRust, align=center]{return: ``I need\\$\mat{A}\vv$ for this $\vv$''};
  % the operator-apply that Palace runs to service the request
  \node[stage, below=12mm of palace, minimum width=30mm, align=center] (apply)
    {\scriptsize apply $\mat{A}=(\mat{K}{-}\sigma\mat{M})^{-1}\mat{M}$\\\scriptsize (the inner stack)};
  \draw[flow] (palace.south) -- (apply.north);
  \draw[flow] (apply.east) -| ([xshift=-3mm]lib.south) node[pos=0.25,below,font=\scriptsize]{$\mat{A}\vv$ back};
  \node[font=\scriptsize\itshape, text=simGray, align=center, anchor=north] at (lib.south|-apply.south)
    {the boundary is a\\\emph{packaging} detail,\\not an algorithm};
\end{tikzpicture}
\caption[The reverse-communication library boundary]{What the ``opaque library call'' of
\cref{ch:eigenvalues} actually is. The Krylov--Schur algorithm of this chapter lives inside the
library (SLEPc or ARPACK, violet box on the right)---but the library cannot call Palace's operator
directly, because it knows nothing of Palace's data structures. So control inverts by
\emph{reverse communication}: Palace calls the library driver \code{naupd} to ``take a step''; the
library advances the eigen-iteration and \emph{returns}, requesting a specific matrix-vector product
$\mat{A}\vv$; Palace runs the entire inner shift-invert stack (\cref{fig:eigenstack}) to compute it
and calls back; the library resumes. The algorithm is the loop of \cref{fig:ksloop}; the only thing
the library boundary adds is this request-and-resume packaging. Open the box and the
``opaque'' eigensolver is exactly the construction of this chapter.}
\label{fig:reversecomm}
\end{figure}

\Cref{fig:reversecomm} draws the dance. The lesson is the one this book returns to again and again:
a black box is opaque only until you build its contents. The eigensolver looked, in
\cref{ch:eigenvalues}, like the one stage of the simulator that was irreducibly someone else's
code---a single library call with no loop of our own. It is not. It is the Rayleigh--Ritz
projection of \cref{sec:rayleighritz}, restarted by the thick-restart of \cref{sec:thickrestart},
built on the Arnoldi and Lanczos recurrences of \cref{ch:orthog}, with the shift-invert linear
solve of \cref{ch:krylov} inside each step. We could write it ourselves; Palace calls the library
because the library is well-tested and tuned, not because the algorithm is beyond us. The box is
demystified.

\begin{notationbox}
\textbf{Why call the library at all, then?} If the algorithm is just this chapter's loop, why not
write it? Three reasons, none algorithmic. First, \emph{robustness}: SLEPc and ARPACK encode
decades of hard-won handling of edge cases---breakdowns, clustered eigenvalues, the
reorthogonalization policies of \cref{sec:lanczosspecific}---that a from-scratch implementation
would rediscover the hard way. Second, \emph{the spectral-transform plumbing}: the library manages
the un-transformation $\lambda = \sigma + 1/\theta$, the Higham scaling, and the convergence
bookkeeping across multiple backends behind one interface. Third, \emph{focus}: Palace's value is in
the physics and the assembly, not in re-implementing a numerical eigensolver. Calling the library is
the right engineering choice precisely \emph{because} we understand what it does---we can configure
it, diagnose it, and trust it, having seen its insides.
\end{notationbox}

\summarysection
The opaque eigenvalue solve of \cref{ch:eigenvalues} is the Krylov--Schur algorithm, and this
chapter built it. \emph{Rayleigh--Ritz projection} takes an orthonormal Krylov basis $\mat{V}_m$
(from the Arnoldi or Lanczos recurrence of \cref{ch:orthog}), forms the small projected operator
$\mat{T}_m = \mat{V}_m^{\!\top}\mat{A}\mat{V}_m$---which is exactly the Hessenberg/tridiagonal matrix
the recurrence already filled---and reads its eigenpairs as \emph{Ritz pairs} $(\theta_i,
\mat{V}_m\vs_i)$ approximating $\mat{A}$'s eigenpairs, the extremal ones first and best (the
polynomial / min-max view). The Ritz residual is the cheap product $\abs{h_{m+1,m}}
\abs{\ve_m^{\!\top}\vs_i}$, a convergence test costing two already-computed scalars. But the basis
cannot grow without bound---storage and orthogonalization cost climb, and Lanczos loses
orthogonality and spawns ghost eigenvalues---so the method must \emph{restart}, and a naive restart
throws away the converged Ritz directions. \emph{Thick-restart Krylov--Schur} (Stewart) is the
elegant fix: reorder the partial Schur form to put the wanted Ritz values first, truncate the basis
to a small size $k$ that retains all the good directions, carry one spill-over vector, and resume
Arnoldi from there---expand, extract, lock, compress, repeat. For symmetric problems the basis
extension is the Lanczos three-term recurrence ($O(1)$ memory, two stored vectors), economical but
demanding a reorthogonalization policy to suppress ghosts. In the eigenmode driver the operator
iterated against is the shift-inverted $(\mat{K} - \sigma\mat{M})^{-1}\mat{M}$ of
\cref{ch:eigenvalues}, so every basis-extension step contains an inner Krylov solve---the full stack
is eigen-iteration $\triangleright$ shift-invert apply $\triangleright$ inner \code{ksp\_solve}
$\triangleright$ preconditioner $\triangleright$ matrix-free operator. SLEPc's \code{EPSKRYLOVSCHUR}
and ARPACK's reverse-communication \code{naupd} run exactly this loop; the ``opaque library call''
is a packaging detail around an algorithm we can---and now do---understand fully.

\furtherreading
The thick-restart Krylov--Schur method is due to Stewart, whose paper~\citep{stewart2001krylovschur}
introduces the Krylov--Schur relation and the reordering-and-truncation restart that
\cref{sec:thickrestart} develops; it is the cleanest single source for the algorithm at the heart of
this chapter, and is the reconstruction the source specification follows. For the broader theory of
Krylov eigensolvers---Rayleigh--Ritz projection, the convergence of Ritz values, the Lanczos and
Arnoldi recurrences, reorthogonalization policies, and the ghost-eigenvalue phenomenon---Saad's
\emph{Numerical Methods for Large Eigenvalue Problems}~\citep{saad2011eig} is the standard graduate
reference and the natural next step after this chapter. The dense QR algorithm that diagonalizes the
small projected $\mat{T}_m$, together with the Schur decomposition and its reordering that
thick-restart relies on, are developed definitively in Golub and Van Loan's \emph{Matrix
Computations}~\citep{golubvanloan2013}, which also treats the symmetric Lanczos process and its
finite-precision behaviour in full. Read Stewart for the restart, Saad for the eigensolver theory,
and Golub and Van Loan for the dense kernels underneath---and keep \cref{ch:eigenvalues}'s
shift-invert transform in view throughout, since it is what makes Krylov projection converge to the
interior modes the simulator actually wants.

\exercisessection
\begin{enumerate}[nosep]
  \item \emph{(Rayleigh--Ritz by hand.)} Repeat \cref{wex:rrritz} on
  $\mat{A} = \begin{psmallmatrix} 3 & 1 & 0 & 0\\ 1 & 3 & 1 & 0\\ 0 & 1 & 3 & 1\\ 0 & 0 & 1 & 3
  \end{psmallmatrix}$ with $\vq_1 = (1,0,0,0)^{\!\top}$. Build $\mat{T}_3$ by Lanczos, compute its
  three Ritz values, and verify that the extreme Ritz values bracket the true extreme eigenvalues
  $3 + 2\cos(k\pi/5)$ from inside.
  \item \emph{(The cheap residual.)} For the largest Ritz pair of \cref{wex:rrritz}, the residual
  was $\abs{h_{4,3}}\abs{\ve_3^{\!\top}\vs} = 1\cdot\tfrac12$. Derive the formula
  $\norm{\mat{A}\vx_i - \theta_i\vx_i} = \abs{h_{m+1,m}}\abs{\ve_m^{\!\top}\vs_i}$ from the Arnoldi
  relation $\mat{A}\mat{V}_m = \mat{V}_{m+1}\bar{\mat{H}}_m$, explaining at each step why the
  projected part cancels. Why does a small $h_{m+1,m}$ make \emph{every} Ritz pair converge at once?
  \item \emph{(Why the extremes.)} Using the polynomial picture $\mathcal{K}_m = \{p(\mat{A})\vv :
  \deg p < m\}$, explain why a degree-$(m{-}1)$ polynomial can isolate an extreme eigenvalue with
  small $m$ but needs large $m$ to isolate an interior one. What does this predict about the relative
  convergence speed of the extreme versus interior Ritz values in \cref{fig:ritzconverge}?
  \item \emph{(The cost of not restarting.)} A basis of size $m$ costs $O(mN)$ storage and
  $O(m^2 N)$ total orthogonalization work for general Arnoldi. If an eigenproblem needs $m = 500$
  steps to converge $10$ modes without restart, but thick-restart caps the basis at $k_{\max} = 50$,
  estimate the ratio of storage and of orthogonalization cost between the two. Which of the three
  pressures of \cref{fig:restartproblem} does restart relieve, and which does reorthogonalization
  relieve?
  \item \emph{(One thick-restart step.)} In \cref{wex:thickrestart} we kept $k = 3$ of $k_{\max} = 6$
  Ritz directions. Suppose instead we want $4$ modes and restart to $k = 5$. Given Ritz values
  $\{7.9, 7.4, 6.8, 6.1, 5.0, 1.2\}$ with residuals $\{10^{-9}, 10^{-7}, 10^{-3}, 10^{-2}, 0.4,
  0.9\}$ and tolerance $10^{-8}$, identify which are locked, which are kept, and which are discarded,
  and explain why keeping one extra direction beyond the wanted count is worthwhile.
  \item \emph{(Lanczos ghosts.)} Explain, in terms of the three-term recurrence's skipped
  projections (\cref{sec:lanczosspecific}), why bare Lanczos produces ghost eigenvalues but bare
  Arnoldi does not. Which reorthogonalization policy of \cref{tab:reorthog} restores safety at the
  least cost, and why does thick-restart partially substitute for reorthogonalization?
  \item \emph{(The shift-invert stack.)} Trace one application of $\mat{A} = (\mat{K} -
  \sigma\mat{M})^{-1}\mat{M}$ to a vector $\vv$ through all five layers of \cref{fig:eigenstack},
  naming the chapter that owns each layer. Why is an eigenmode analysis's runtime dominated by the
  inner \code{ksp\_solve} rather than by the outer eigen-iteration? What does this imply about the
  payoff of a better preconditioner (\cref{ch:multigrid}) for the eigenmode driver specifically?
  \item \emph{(Reverse communication.)} Explain why ARPACK uses a reverse-communication interface
  rather than calling Palace's operator directly, and describe the request-and-resume cycle of
  \cref{fig:reversecomm}. In what sense is the library boundary a ``packaging detail, not an
  algorithm''? Give one concrete advantage of calling the library over writing the Krylov--Schur loop
  from scratch, and one concrete advantage of \emph{understanding} the loop even though you call the
  library.
\end{enumerate}
